\documentclass[10pt,a4paper]{amsart}
\usepackage[margin=19mm]{geometry}
\usepackage{amsmath,amssymb,mathtools}
\usepackage[scr=rsfs,cal=euler]{mathalfa}
\usepackage{xfrac}
\usepackage[unicode,hidelinks,bookmarks=false]{hyperref}
\newcommand{\ie}{i.e.\ }
\newcommand{\cf}{cf.\ }
\newcommand{\resp}{respectively\ }
\newcommand{\Subsection}{\subsection}
\providecommand{\nobreakdash}{\nobreak\hbox{-}}
\setlength{\emergencystretch}{2em}
\multlinegap0pt
\usepackage{amssymb}
\usepackage{natbib}
\usepackage{xcolor}
\newtheorem{lemma}{Lemma}
\usepackage{cleveref}
\crefname{lemma}{Lemma}{Lemmas}
\title{Long cycles in vertex transitive digraphs}
\author{Matija Buci\'c, Kevin Hendrey, Bojan Mohar, Raphael Steiner, Liana Yepremyan}
\date{Source: arXiv:2602.16333v1 (CC BY 4.0)}
\begin{document}
\maketitle

\noindent\emph{Source and licence.} Long cycles in vertex transitive digraphs. Version arXiv:2602.16333v1 (CC BY 4.0), distributed by arXiv under the Creative Commons Attribution 4.0 International licence, http://creativecommons.org/licenses/by/4.0/, as recorded in the arXiv licence metadata for this exact version and shown on the abstract page https://arxiv.org/abs/2602.16333v1. Full licence notice: \texttt{licenses/arxiv-2602.16333-CC-BY-4.0.txt}.\par\smallskip

\noindent\textit{Editorial scope.} Counted unit: the article's lemma lem:long-cycle-in-cycle-graphs (source0.tex lines 307--309). The article's complete original proof of it is delivered with it (source0.tex lines 310--350, one \texttt{proof} environment of 41 lines). No mathematics is written, repaired or re-derived: the statement and the proof are the article's own lines, in the article's own notation, and no retained line is substituted or re-flowed.  No further source-local context is needed: every cross-reference of every retained line resolves inside the two delivered spans.  One declared adaptation: exactly 1 ancillary strings inside the retained spans are omitted (image-inclusion commands and, where the source repeats a label, the repeated label definitions) and no other line differs from the source; the figure environments, with their placement, captions and labels, are kept, and the package records the omitted strings in fidelity receipts bound to the affected bodies.  No retained line contains a citation, so the wrapper carries no bibliography and no bibliography entry.  The retained lines contain the article's own diagram code, which the wrapper typesets with the standard tikz packages; no image file is delivered and the package declares no asset dependency.  Editorial formatting only, in addition: an AMS \texttt{amsart} preamble that restates the article's own declarations and macros used by the retained lines, namely the environments \texttt{lemma} on separate counters, together with the standard packages those lines need.  The retained environments print in this document as Lemma 1 in order of appearance, whereas the article prints the same items with the numbering of its own full document.  The complete native source of the article as distributed in the arXiv e-print for this exact version is bundled unchanged as \texttt{sources/194872/source0.tex}.
\begin{lemma}\label{lem:long-cycle-in-cycle-graphs}
	In any connected, nearly transitive graph $G$ with diameter $d\ge 20$, there is an induced cycle of length at least $d-17$.
\end{lemma}

\begin{proof}
	Suppose towards a contradiction that every induced cycle in $G$ has length less than $\frac{d}{4}$. 

	Let us fix a shortest path $S$ between two vertices $v,u$ at distance $d$ in $G$. Note that $S$ is a geodesic path and is hence induced. 
	Let $m$ be a vertex of $S$ at distance at least $\frac{d-1}{2}$ from both $u$ and $v$. 
	Let $L$ be the subpath of $S$ between $v$ and $m$ and let $R$ be the subpath between $m$ and $u$.
	
	Let $P$ be an induced path in $G$ of maximum length subject to the condition that the subpath $Q$ of $P$ on its last $\lceil \frac{d-5}{2}\rceil$ vertices is a geodesic path. Note that $S$ satisfies this condition and so such a path exists and has length at least $d$. Let $w$ be the common endpoint of $Q$ and $P$, let $x$ be the other endpoint of $P$ and let $y$ be the other endpoint of $Q$.
	
	Let $S'$ be the image of $S$ under an automorphism $\phi$ mapping $m$ to $w$ or a neighbor of $w$. Let %$v'=\phi(v), u'=\phi(u),$ 
    $L':=\phi(L), R':=\phi(R),$ $v':=\phi(v)$, $u':=\phi(u)$ and $w':=\phi(m)$.

    We first claim that either $L'$ or $R'$ intersects $Q \cup N(Q)$ only in vertices at distance at most $3$ from $w'$. To see this, let $a'$ and $b'$ be vertices of $L' \cap (Q \cup N(Q))$ and $R' \cap (Q \cup N(Q))$ respectively. Let $a$ and $b$ be vertices of $Q$ at distance at most one from $a'$ and $b'$, respectively. See \Cref{fig:translation} for an illustration of the setup so far.

    \begin{figure}[h]
    \centering

    \caption{Illustration of the setup in the proof of \Cref{lem:long-cycle-in-cycle-graphs}. Dashed lines depict pairs of vertices at distance up to one. $L'\cup R'=S'$ is a geodesic path, and so is $Q$. This implies that $d(w,a)$ is very close to $d(w,a')$ and $d(w,b)$ to $d(w,b')$. The argument now relies on the fact that $d(a',b')$ is roughly the sum of these two distances since $a',b'$ are both vertices of the geodesic path $S'$, but also by going through the $a$-$b$ segment, we can find a path of distance roughly the difference between these two distances. The conclusion is that one of these distances needs to be very small.}
    \label{fig:translation}
\end{figure}

    Note that since $d(w,w'),d(a,a'),d(b,b') \le 1$, we get via triangle inequality that 
    \begin{align}\label{eq:1}
        |d(a',w')-d(a,w)|&\le d(a',a)+d(w,w') \le 2 \quad \text{ and}\\
        \label{eq:2}
        |d(b',w')-d(b,w)|&\le d(b',b)+d(w,w') \hspace{0.05cm}\le 2.
    \end{align}
    
    Since $S'$ is geodesic (being an image of a geodesic path under an automorphism), so are all of its subpaths. Combining this with several applications of the triangle inequality, we get
    \begin{equation}\label{eq:3}
        d(a', w')+d(b',w')=d(a',b') \le d(a,b)+2 = |d(a,w)-d(b,w)|+2.
    \end{equation} 
    %Similarly, since $Q$ is geodesic we have $d(x,w)\le d(a,w)+2$ and $d(y,w) \le d(b,w)+2$. 
    Combining \eqref{eq:1}, \eqref{eq:2} and \eqref{eq:3} we obtain 
	$$d(a',w')+d(b',w') \le |d(a',w')-d(b',w')|+6 \implies \min\{d(a',w'),d(b',w')\}\le 3,$$
	 as claimed.
	
	Suppose without loss of generality
    that $L'$ intersects $Q \cup N(Q)$ only in vertices at distance at most $3$ from $w'$. Pick a vertex $c$ in this intersection farthest from $w'$. Let $z$ be a vertex at distance at most one from $c$ on $Q$, farthest from $w$. By the triangle inequality, we know that $d(z,w) \le d(c,w') +2 \le 5$. Now, let $Q'$ be the subpath of $Q$ from $y$ to $z$, and let $L''$ be the subpath of $L'$ from $c$ to $v'$. Note that $V(Q') \cup V(L'')$ induces a path by construction (since $c$ and $z$ are equal or adjacent). Let $P'$ denote the subpath of $P$ from $x$ to $y$.
    We now claim that there must be a pair of vertices $s\in V(P')$ and $t\in V(L'')$ with $d(s,t)\le 1$. Indeed, suppose for a contradiction such a pair does not exist. Then $P'$ and $L''$ are vertex-disjoint and have no connecting edges. In particular, appending the path induced by $V(Q')\cup V(L'')$ to $P'$, we obtain an induced path whose last $\lceil \frac{d-5}{2}\rceil$ vertices induce a subpath of $L'$ and thus a geodesic path. Moreover, the total length of the new path is at least $|E(P)|+\frac{d-1}{2}-8+1 > |E(P)|$, and this contradicts our choice of $P$. Hence, there indeed need to exist vertices $s\in V(P'), t \in V(L')$ with $s=t$ or $st\in E(G)$. Pick such $s$ and $t$ for which $d_{P'}(s,y)+d_{L''}(t,c)$ is minimized. Let $C$ be the cycle formed by the union of the segment of $P'$ from $s$ to $y$, the segment of $L''$ from $t$ to $c$, the edge $st$ if $s\neq t$, as well as the path in $G$ induced by $V(Q')\cup V(L'')$. It is not hard to see that by our choice of $c,z,s,t$ respectively $C$ forms an induced cycle in $G$. Since $C$ by definition contains $Q'$ as a subpath and since $Q'$ is geodesic, we find that $|E(C)|\ge 2|E(Q')|\geq 2(\frac{d-7}{2}-5)=d-17$, as desired. 
\end{proof}

\end{document}
